\section{Exact endpoint moments and polynomial Bregman asymmetry}
\label{sec:p9-exact-asymmetry}

The endpoint first-moment problem can be solved exactly by positive
quadrature. Its sum-of-squares version is already the linear-objective
case in de Klerk--Laurent, \emph{Worst-case examples for Lasserre's
measure--based hierarchy for polynomial optimization on the hypercube},
\cite[Theorem~3.1 and Corollary~3.2]{deklerk}. The reduction below is a
classical quadrature consequence, not a new approximation principle.
It strengthens the constants and identifies extremizers for one ingredient;
it does not determine a hierarchy price.

For an integer $r\ge0$, define
\[
 \lambda_r=\inf_{\substack{g\not\equiv0,\ \deg g\le r\\
                                      g\ge0\text{ on }[0,1]}}
       \frac{\int_0^1 t g(t)\,dt}{\int_0^1g(t)\,dt}.
\]
For $\alpha,\beta>-1$, write $\xi_j^{(\alpha,\beta)}$ for the
smallest zero of the degree-$j$ Jacobi polynomial
$P_j^{(\alpha,\beta)}$ on $[-1,1]$, with weight
$(1-x)^\alpha(1+x)^\beta$.

\begin{theorem}[Exact endpoint moment]
For every $m\ge0$,
\begin{equation}
 \lambda_{2m}=\frac{1+\xi_{m+1}^{(0,0)}}2,
 \qquad
 \lambda_{2m+1}=\frac{1+\xi_{m+1}^{(1,0)}}2.
 \label{eq:p9-exact-lambda}
\end{equation}
The infima are attained. In particular,
$\lambda_r=\Theta((r+1)^{-2})$.
\end{theorem}
\begin{proof}
For even degree, take the $m+1$-node Gauss--Legendre rule on $[0,1]$
\cite[Section~3.5(v)]{dlmf}.
Its nodes $t_1<\cdots<t_{m+1}$ are the shifted Legendre zeros, its
weights $\omega_j$ are positive, and it integrates polynomials of degree
at most $2m+1$ exactly. Thus for every admissible $g$,
\[
 \int_0^1 t g(t)\,dt-t_1\int_0^1g(t)\,dt
  =\sum_{j=1}^{m+1}\omega_j(t_j-t_1)g(t_j)\ge0.
\]
The polynomial $g(t)=\prod_{j=2}^{m+1}(t-t_j)^2$ has degree $2m$,
is nonnegative, and makes the right side zero. The empty product for
$m=0$ is one. This proves the first formula.

For clarity, the requisite odd-degree quadrature rule is constructed here.
Let $J$ be the degree-$m+1$ polynomial orthogonal for weight $1-t$ on
$[0,1]$. Its $m+1$ simple interior zeros are $t_1<\cdots<t_{m+1}$;
add the node $t_{m+2}=1$. Integrate the Lagrange interpolant on these
$m+2$ nodes to define an interpolatory rule. For a polynomial $f$ of
degree at most $2m+2$, divide
\[
 f(t)=(t-1)J(t)q(t)+s(t),\qquad
 \deg q\le m,\quad\deg s\le m+1.
\]
The first summand integrates to zero by orthogonality and vanishes at all
nodes. The rule therefore integrates $f$ exactly. If $L_j$ is the
degree-$m+1$ cardinal Lagrange polynomial, exactness for $L_j^2$ gives
$\omega_j=\int_0^1L_j(t)^2\,dt>0$.
Applying this positive rule to $g$ and $tg$ proves that the moment ratio
is at least $t_1$. Equality is attained by
\[
 g(t)=(1-t)\prod_{j=2}^{m+1}(t-t_j)^2,
\]
of degree $2m+1$. The orthogonality weight identifies $J$ with the
shifted $P_{m+1}^{(1,0)}$, proving the second formula.

For both fixed parameter pairs, the smallest Jacobi zero is
$-1+\Theta((m+1)^{-2})$. This standard extreme-zero estimate is recorded
in de Klerk--Laurent~\cite[Corollary~2.3]{deklerk}, and gives the stated order.
\end{proof}

\begin{theorem}[Sharp oriented polynomial asymmetry]
Let $\varphi''=g$ be a nonnegative nonzero polynomial of degree at most
$r$ on an interval $[u,v]$, $u<v$. Then
\begin{equation}
 \max\left\{\frac{B_\varphi(v,u)}{B_\varphi(u,v)},
             \frac{B_\varphi(u,v)}{B_\varphi(v,u)}\right\}
 \le\frac{1-\lambda_r}{\lambda_r}.
 \label{eq:p9-exact-asymmetry}
\end{equation}
The bound is attained within nonnegative polynomial curvature, is the
supremum if curvature is required to be strictly positive everywhere on
the interval, and is $\Theta((r+1)^2)$.
\end{theorem}
\begin{proof}
Affine rescaling reduces to $[0,1]$. Writing
$\mu=\int_0^1g$ and $a=\mu^{-1}\int_0^1tg$ gives
\[
 B_\varphi(0,1)=a\mu,
 \qquad B_\varphi(1,0)=(1-a)\mu.
\]
The preceding theorem and its reflected version imply
$\lambda_r\le a\le1-\lambda_r$. The ratio bound follows, and the
extremizing curvatures above attain equality. Such curvature still
induces strict convexity because it vanishes on no nonempty interval.
For everywhere-positive curvature, add $\varepsilon>0$ and send it to
zero; all displayed integrals converge and the denominator stays positive.
The extreme-zero order proves the last assertion.
\end{proof}

\paragraph{Why oriented asymmetry is not a hierarchy lower bound.}
After normalizing $\int_0^1g=1$, a first-moment estimate
$\int_0^1tg=O(r^{-2})$ gives only
$\int_d^1g\le O((r^2d)^{-1})$. At $d\asymp r^{-2}$ this is a
constant bound, whereas the existing four-state endpoint proof uses a
leakage bound much smaller than a power of $d$. The exact moment theorem
does not supply that stronger assertion. Conversely, sharp oriented
asymmetry need not generate incompatible optimal partitions at two
capacities. It cannot be substituted for a hierarchy-price lower bound.
The general-source $O(r^2)$ two-capacity bound in Phase~8 uses an anchor
quasi-triangle inequality, a different optimization from the endpoint
orientation ratio solved here. The sharp four-state result in
Theorem~\ref{thm:weighted-sharp} instead uses a direct kernel argument.
